\chapter{绪论}

本章的主要内容与学校提供的Word模板中内容一致，图片与表格均采用原始设定大小，%
主要是为了说明格式的统一。%
但是，\LaTeX{}的一些禁则，专业排版的能力，对公式及文献的处理都是得天独厚的，%
我们不必刻意去追求与Word的完美匹配。而且你将会发现，用\LaTeX{}书写论文的美！ %

\section{（1.1 题目）}
正文内容

\subsection{（1.1.1 题目）}
正文内容

正文内容

\begin{figure}[htp]
	\centering
	\includegraphics{picmain}
	\caption{图 1.1 名称}
\end{figure}

\subsubsection{（1.1.1.1 题目）}
正文内容

正文内容

正文内容

\subsubsection{（1.1.1.2 题目）}
正文内容

正文内容

正文内容

\subsection{（1.1.2 题目）}
正文内容

正文内容

\begin{figure}[htp]
	\centering
	\includegraphics{picmain}
	\caption{图 1.2 名称}
\end{figure}

\section{（1.2 题目）}
正文内容

正文内容

\begin{table}[htp]
	\centering
	\begin{minipage}[t]{0.8\linewidth} % 如果想在表格中使用脚注，minipage是个不错的办法
		\caption[表 1.1 名称]{}
		\begin{tabular*}{\textwidth}{lp{10cm}}
			\toprule[1.5pt]
			{\hei 列1} & {\hei 列2} \\
			\midrule[1pt]
			&  \\
			& \\
			& \\
			& \\
			& \\
			& \\
			\bottomrule[1.5pt]
		\end{tabular*}
	\end{minipage}
\end{table}

正文内容

正文内容

正文内容

正文内容

\section{（1.3 题目）}
正文内容

正文内容

正文内容

正文内容

正文内容

正文内容

\subsection{（1.3.1 题目）}
正文内容

\begin{figure}[htp]
	\centering
	\includegraphics{picmain}
	\caption{图 1.3 名称}
\end{figure}

\subsection{（1.3.2 题目）}
正文内容

正文内容

\begin{table}[htp]
	\centering
	\begin{minipage}[t]{0.8\linewidth} % 如果想在表格中使用脚注，minipage是个不错的办法
		\caption[表 1.2 名称]{}
		\begin{tabular*}{\textwidth}{lp{10cm}}
			\toprule[1.5pt]
			{\hei 列1} & {\hei 列2} \\
			\midrule[1pt]
			&  \\
			& \\
			& \\
			& \\
			& \\
			& \\
			\bottomrule[1.5pt]
		\end{tabular*}
	\end{minipage}
\end{table}

